\documentclass[10pt,letterpaper]{article}

% --- Paquetes de tipografía, color y maquetación ---
\usepackage[utf8]{inputenc}
\usepackage[spanish,es-nodecimaldot]{babel}
\usepackage[margin=1.5cm,top=1.3cm,bottom=1.3cm]{geometry}
\usepackage{amsmath,amssymb}
\usepackage{booktabs}
\usepackage{enumitem}
\usepackage{xcolor}
\usepackage{hyperref}
\usepackage{titlesec}
\usepackage{microtype}
\usepackage{tabularx}
\usepackage{tcolorbox}

% --- Configuración de tcolorbox ---
\tcbuselibrary{skins,breakable}

% --- Paleta cromática profesional y armónica ---
\definecolor{primary}{RGB}{20, 45, 80}       % Azul marino institucional
\definecolor{secondary}{RGB}{35, 110, 160}   % Azul medio
\definecolor{bgsoft}{RGB}{244, 248, 253}     % Fondo azul suave
\definecolor{cardbg}{RGB}{250, 251, 252}     % Fondo gris claro
\definecolor{cardborder}{RGB}{210, 218, 228} % Borde sutil
\definecolor{accenturl}{RGB}{13, 94, 158}    % Azul enlace

% --- Configuración de hipervínculos ---
\hypersetup{
    colorlinks=true,
    linkcolor=accenturl,
    urlcolor=accenturl,
    citecolor=accenturl
}

% --- Estilos de Título de Secciones ---
\titleformat{\section}
  {\large\bfseries\color{primary}}
  {\thesection.}{0.4em}{}
  [{\color{primary}\titlerule[1.1pt]}]

\titleformat{\subsection}
  {\normalsize\bfseries\color{secondary}}
  {\thesubsection.}{0.35em}{}

\titlespacing*{\section}{0pt}{6pt}{3pt}
\titlespacing*{\subsection}{0pt}{4pt}{2pt}

% --- Espaciado de párrafos y listas ---
\setlength{\parskip}{2.8pt}
\setlength{\parindent}{0pt}
\setlist[itemize]{leftmargin=14pt, itemsep=2pt, topsep=1.5pt, parsep=0pt}
\setlist[enumerate]{leftmargin=14pt, itemsep=2pt, topsep=1.5pt, parsep=0pt}

% --- Cajas de diseño estructurado ---
\newtcolorbox{highlightbox}[1][]{
    colback=bgsoft,
    colframe=secondary,
    arc=2.5mm,
    boxrule=0.8pt,
    left=7pt, right=7pt, top=4pt, bottom=4pt,
    #1
}

\newtcolorbox{datasetbox}[2][]{
    colback=cardbg,
    colframe=cardborder,
    arc=2mm,
    boxrule=0.6pt,
    left=6pt, right=6pt, top=3pt, bottom=3pt,
    title={\footnotesize\bfseries\color{primary}#2},
    #1
}

\begin{document}

% =========================================================================
% ENCABEZADO INSTITUCIONAL
% =========================================================================
\begin{tcolorbox}[colback=primary, colframe=primary, arc=2mm, top=5pt, bottom=5pt, left=8pt, right=8pt]
    \centering
    {\fontsize{10}{12}\selectfont\textbf{\color{white}MAESTRÍA EN INTELIGENCIA ARTIFICIAL -- MATEMÁTICAS Y PROGRAMACIÓN IA}}\\[1.5pt]
    {\fontsize{13}{15}\selectfont\textbf{\color{white}PROPUESTA DE PROYECTO DE INVESTIGACIÓN Y DESARROLLO}}
\end{tcolorbox}

\vspace{-4pt}

% =========================================================================
\section{Tema}
% =========================================================================
\begin{tcolorbox}[colback=bgsoft, colframe=secondary, arc=2mm, boxrule=0.9pt, left=7pt, right=7pt, top=3.5pt, bottom=3.5pt]
    \centering
    \textbf{\color{primary}Detección de suplantación de identidad frente a variaciones naturales del comportamiento mediante biometría conductual, Feature Engineering y Machine Learning}
\end{tcolorbox}

% =========================================================================
\section{Descripción de la Propuesta}
% =========================================================================

\subsection{Descripción Funcional}
Los sistemas tradicionales de autenticación estática (contraseñas, tokens) verifican la identidad únicamente al inicio de la sesión. Frente a esta limitación, la \textbf{autenticación continua} mediante \textbf{biometría conductual} evalúa en segundo plano los patrones neuromotores de interacción del usuario (dinámica de tecleo y movimientos del cursor) de forma transparente y sin fricción.

\textbf{Problema central:} El comportamiento de una persona no es constante en el tiempo; experimenta \textbf{variaciones naturales intra-usuario} (\textit{behavioral drift}) ocasionadas por fatiga, estrés o cambios de contexto. Si el sistema no reconoce esta variabilidad inherente, genera altos índices de falsas alarmas (\textit{False Rejection Rate}). 

\begin{highlightbox}[title=\footnotesize\textbf{Pregunta Central de Investigación}]
    \textit{¿Qué características espaciotemporales extraídas mediante Feature Engineering permiten discriminar con robustez entre la deriva natural del usuario legítimo y una suplantación real de identidad, minimizando los falsos rechazos?}
\end{highlightbox}

El sistema analizará y diferenciará entre tres estados operativos:
\begin{enumerate}
    \item \textbf{Comportamiento habitual:} Patrón de interacción recurrente del usuario legítimo.
    \item \textbf{Variaciones naturales:} Desviaciones legítimas de velocidad y pausas del mismo usuario entre distintas sesiones.
    \item \textbf{Suplantación de identidad:} Patrón anómalo correspondiente a un atacante o usuario no autorizado.
\end{enumerate}

\textbf{Datasets Públicos y Gratuitos de Libre Descarga:}
\begin{itemize}
    \item \textbf{Balabit Mouse Dynamics Challenge Dataset:} Benchmark público internacional con sesiones continuas de interacción de múltiples usuarios en entornos reales, capturando coordenadas $(x, y, t)$, clics y trayectorias.\\
    \textbf{Enlace de descarga:} \url{https://github.com/balabit/Mouse-Dynamics-Challenge}
    \item \textbf{Keystroke and Mouse Dynamics for UEBA Dataset (Mendeley Data):} Repositorio multi-sesión abierto con más de 140\,000 registros integrados de pulsaciones de teclado (\textit{Hold Time, Flight Time}) y métricas de cursor.\\
    \textbf{Enlace de descarga:} \url{https://data.mendeley.com/datasets/f78jsh6zp9/2}
    \item \textbf{CMU Keystroke Dynamics Benchmark Dataset:} Dataset de referencia histórica de dinámica de tecleo multi-sesión.\\
    \textbf{Enlace de descarga:} \url{https://www.cs.cmu.edu/~keystroke/}
\end{itemize}

\textbf{Estructuración de Clases y Escenario de Clasificación:}
\begin{itemize}
    \item \textbf{Clasificación Binaria Supervisada:} \textbf{Clase 0 (Legítimo)}, agrupando sesiones base y sesiones con variación natural del usuario; frente a \textbf{Clase 1 (Impostor)}, con sesiones de terceros.
    \item \textbf{Detección de Anomalías (One-Class):} Perfilado no supervisado entrenado exclusivamente con el comportamiento del usuario legítimo para evaluar umbrales dinámicos frente a ataques.
\end{itemize}

\newpage

\subsection{Descripción Técnica}
El proyecto se implementará en \textbf{Python}, empleando las siguientes librerías especializadas:
\begin{itemize}
    \item \textbf{Procesamiento Numérico y Datos:} \texttt{pandas}, \texttt{numpy}, \texttt{scipy} (álgebra lineal, transformaciones estadísticas y series temporales).
    \item \textbf{Machine Learning \& Validación:} \texttt{scikit-learn} (clasificadores supervisados, reducción dimensional y métricas).
    \item \textbf{Interpretabilidad y Visualización:} \texttt{matplotlib}, \texttt{seaborn} y \texttt{SHAP} (valores de Shapley para análisis de atribución).
\end{itemize}

\textbf{Pipeline y Arquitectura de la Solución:}
\begin{enumerate}
    \item \textbf{Ingesta y Preprocesamiento:} Limpieza de valores nulos, filtrado de eventos espurios y segmentación en ventanas temporales móviles por usuario y sesión.
    \item \textbf{Feature Engineering de Alto Nivel:} Extracción de descriptores conductuales especializados:
    \begin{itemize}
        \item \textit{Teclado:} Tiempos de pulsación (\textit{Hold Time}), tiempos de vuelo entre dígrafos (\textit{Flight Time}), velocidad de tipeo (WPM), variabilidad del ritmo, frecuencia de borrado (\textit{Backspace}) y duración de pausas cognitivas.
        \item \textit{Mouse:} Velocidad tangencial media, aceleración instantánea, curvatura y rectitud de trayectorias (distancia euclidiana vs. real), tasa de cambio de dirección angular (\textit{jitter}), frecuencia de clics y ratio de tiempo inactivo.
    \end{itemize}
    \item \textbf{Análisis Estadístico y Selección de Atributos:} Cálculo de varianza intra-sesión vs. inter-usuario para seleccionar variables con alta estabilidad temporal individual y alto poder discriminativo frente a terceros.
    \item \textbf{Entrenamiento y Comparación de Modelos:} Se contrastarán dos configuraciones experimentales:
    \begin{itemize}
        \item \textit{Modelo Baseline:} Entrenado únicamente con variables crudas o estadísticas elementales directas.
        \item \textit{Modelo con Feature Engineering:} Entrenado con el conjunto optimizado de descriptores espaciotemporales.
    \end{itemize}
    Modelos evaluados: \textbf{Logistic Regression}, \textbf{Random Forest}, \textbf{Support Vector Machines (SVM RBF)} e \textbf{Isolation Forest}.
    \item \textbf{Validación Estratificada por Sesión (Out-of-Session Split):} Partición estricta basada en sesiones temporales independientes (las primeras sesiones entrenan y sesiones posteriores evalúan), garantizando que el modelo generalice ante la deriva natural sin memorizar una única sesión (\textit{data leakage}).
    \item \textbf{Evaluación e Interpretabilidad:} Medición mediante \textit{Precision, Recall, F1-Score, ROC-AUC, EER} (Equal Error Rate) y análisis de importancia de variables con \texttt{SHAP}.
\end{enumerate}

\begin{tcolorbox}[colback=bgsoft, colframe=secondary, arc=2mm, boxrule=0.7pt, left=7pt, right=7pt, top=3pt, bottom=3pt]
    \small\textbf{\color{primary}Tema de Investigación Complementario (Fuera del Temario del Curso):}\\
    \footnotesize Se investigarán los fundamentos teóricos de la \textbf{Biometría Conductual Continua} y el fenómeno de \textbf{Behavioral Drift} (deriva temporal del patrón de conducta), analizando técnicas de actualización adaptativa de perfiles y ajuste dinámico de umbrales para mantener la seguridad sin generar fricción al usuario legítimo.
\end{tcolorbox}

% =========================================================================
\section{Objetivos}
% =========================================================================

\subsection{Objetivo General}
Desarrollar y evaluar un sistema de Machine Learning fundamentado en Feature Engineering para la autenticación continua mediante biometría conductual de teclado y mouse, capaz de discriminar de forma robusta entre variaciones naturales del comportamiento de un usuario legítimo y eventos de suplantación de identidad.

\subsection{Objetivos Específicos}
\begin{enumerate}
    \item \textbf{Adquirir y estructurar} datasets públicos de biometría conductual (teclado y mouse) con sesiones independientes multi-usuario.
    \item \textbf{Implementar un pipeline de Feature Engineering} que extraiga descriptores cinemáticos, estadísticos y temporales.
    \item \textbf{Evaluar estadísticamente} la estabilidad temporal intra-usuario y el poder discriminativo inter-usuario de las variables.
    \item \textbf{Entrenar y comparar} múltiples modelos de Machine Learning bajo las modalidades \textit{Baseline} y con \textit{Feature Engineering}.
    \item \textbf{Validar la robustez del sistema} ante variaciones naturales aplicando evaluación por sesiones desacopladas (\textit{out-of-session}).
    \item \textbf{Interpretar el rendimiento} mediante métricas de clasificación (\textit{Precision, Recall, F1, ROC-AUC, EER}) y análisis SHAP.
    \item \textbf{Investigar y sintetizar} los conceptos teóricos de biometría conductual, autenticación continua y \textit{behavioral drift}.
\end{enumerate}

\end{document}